\section{Periodic rates and nonnegative full densities}\label{sec:periodic}

The selected first characteristic mode does not by itself produce a
periodic density. We first remove all other harmonic resonances, solve
a scalar periodic rate equation, and finally identify its signed
linear response with a positive prescribed evolution.

\subsection{Physical parameter dependence and harmonic selection}

All parameter perturbations in this section use the original physical
variables on the unchanged half-strip. In particular the floor does
not move when $m$ or $a$ changes.

\begin{lemma}[Analyticity at positive noise]\label{lem:parameter-analyticity}
Locally at every $a>0$ under consideration, the completed inverse
$R_{i\theta}B:\HH_\sigma\to\EE_\sigma$ is holomorphic in $(m,a)$,
uniformly in $\theta\in\R$ on a sufficiently small complex parameter
neighborhood. The stationary triple $(M,M_y,M_F)$ and $N_0=\ell M$
are locally analytic. The response \eqref{eq:physical-response} is
locally analytic in $(\lambda,c,m,a)$ near every imaginary frequency.
The selected crossing functions $\cH,\OmH$ are real analytic on
$U=(0,a_0)\times(2,3)$.
\end{lemma}

\begin{proof}
At a fixed base $(m,a)$, regard the density, derivative, and outgoing
trace as separate components of $\EE_\sigma$. A perturbed inverse solves
\begin{equation}\label{eq:parameter-inverse}
 h=R_{i\theta}B[g-\delta_m h+\delta_a h_y].
\end{equation}
The right side uses the first and second components of $h$. The
uniform triple bound makes its linear perturbation small on the
ambient product for small complex $\delta_m,\delta_a$, independently
of $\theta$. Its Neumann series is analytic. Every output has its
derivative and trace identified by the completed weak equation;
closedness of the distributional derivative and passage in that
equation retain these identities in the limit. Its floor flux is
the perturbed natural flux. For real parameters, homogeneous
uniqueness identifies it with the actual prescribed inverse.

Likewise the nearby normalized stationary density is $M+h$, where
\begin{equation}\label{eq:parameter-stationary}
 h=R_0B[-\delta_m(M+h)+\delta_a(M_y+h_y)],\qquad \int_Dh=0.
\end{equation}
The same series gives analyticity. For real parameters normalized
stationary uniqueness identifies this density with the stationary
state supplied by Proposition~\ref{prop:frozen}, including its
positivity. Weighted norm closeness alone is not used for that
conclusion. Local frequency analyticity follows from the ordinary
resolvent identity as in Proposition~\ref{prop:divergence}. This
proves the response statement.

The simple zero and the nonzero real coupling derivative in
\eqref{eq:crossing} and \eqref{eq:response-root} give analytic crossing
functions by the implicit function argument. Local solutions patch
on $U$ by uniqueness of the selected branch and its crossing.
No uniform radius in $m$ as $a\downarrow0$, and no convergence of
parameter derivatives in that limit, is asserted.
\end{proof}

\begin{lemma}[A base point with only first-harmonic resonance]\label{lem:selection}
There is $(a,m_*)\in U$ such that, writing
\[
 c_*=\cH(a,m_*),\qquad \Omega_*=\OmH(a,m_*),\qquad T=2\pi/\Omega_*,
\]
one has
\begin{equation}\label{eq:nonresonance}
 \partial_m\OmH(a,m_*)\ne0,\qquad
 1-\resp(ik\Omega_*;m_*,a,c_*)\ne0
                         \quad(k\in\mathbb Z\setminus\{-1,1\}).
\end{equation}
The first-harmonic zeros are simple.
\end{lemma}

\begin{proof}
For each fixed $k\ge2$, Proposition~\ref{prop:return-root} gives, at
$c=\cH$,
\[
 \operatorname{Re}\lambda_k
 =a\left[-k^2D+\frac D{\gamma_1}\gamma_k\right]+O(a^{3/2})<0
\]
for small $a$ depending on this $k$. The strict sign follows from
$\gamma_k/k^2<\gamma_1$. The real analytic complex-valued function
\[
 r_k(a,m)=1-\resp(ik\OmH(a,m);m,a,\cH(a,m))
\]
is not identically zero on $U$. Indeed, for this fixed $k$ take $a$
small enough that $ik\OmH$ lies in its fixed return-root disk.
A zero of $r_k$ would, by the imaginary-axis matching in
\eqref{eq:response-match}, give an eigenvalue-one equation for $K$
there. Its only root is $\lambda_k$, which has negative real part,
a contradiction.

Also $\OmH(a,m)\to\omega(m)$ uniformly on $I$, while
\[
 \tau_m=-\int_0^L[m-b(x)]^{-2}\dd x<0.
\]
Choose two interior means with different $\omega$. For all sufficiently
small $a$, $\OmH(a,\cdot)$ has different values at those means.
Thus $\partial_m\OmH$ is not the zero analytic function on $U$.
The zero set of a nonzero real analytic function on this connected
open set has empty interior: vanishing on a ball makes all its
power-series coefficients zero and propagates along a chain of
overlapping balls. For complex-valued $r_k$, apply this observation
to a nonzero real or imaginary part.

The zero sets of all $r_k$ and of $\partial_m\OmH$ are closed with
empty interior in $U$. Successively choose nested closed balls,
each inside the preceding interior, avoiding the next zero set,
with radii tending to zero and the first ball compactly inside $U$.
Completeness gives a point avoiding every set. The static multiplier
is nonzero by \eqref{eq:static-response}; negative harmonics are
complex conjugates. This proves \eqref{eq:nonresonance}.
Simplicity at the first harmonics follows from
\eqref{eq:response-root}.
\end{proof}

The selection is existential and uses only a fixed-$k$ asymptotic
for each integer $k$. It does not require a small-noise expansion
uniform over infinitely many harmonics.

\subsection{The periodic inverse and the rate equation}

From now on fix the point and the period in Lemma~\ref{lem:selection}.
Use Fourier coefficients
$u_k=T^{-1}\int_0^T u(t)e^{-ik\Omega_*t}\dd t$, so that
\[
 \norm u_{H^1(0,T;E)}^2
   =T\sum_{k\in\mathbb Z}(1+k^2\Omega_*^2)\norm{u_k}_E^2
\]
for a Hilbert space $E$.

\begin{lemma}[High-frequency response and periodic inverse]\label{lem:periodic-inverse}
At fixed positive parameters,
\begin{equation}\label{eq:response-tail}
 \resp(i\theta;m,a,c)\longrightarrow0\quad(|\theta|\to\infty).
\end{equation}
The coefficientwise map
\[
 (\mathcal R_TBg)_k=\mathcal R_{k\Omega_*}g_k
\]
is bounded from $H^1_{\per}(0,T;\HH_\sigma)$ to
$H^1_{\per}(0,T;\EE_\sigma)$, with its density, gradient, trace, and
floor identities. The linear rate-response multiplier on
$H^1_{\per}(0,T)$ is compact.
\end{lemma}

\begin{proof}
For fixed $h\in\HH_{\sigma,0}$, $t\mapsto S(t)h$ is Bochner
integrable by the frozen gap. Approximate it in $L^1_t\HH_\sigma$
by finitely many step functions and integrate their scalar
exponentials. Formula \eqref{eq:ordinary-resolvent} then shows
$R_{i\theta}h\to0$ in $\HH_\sigma$. Apply
\eqref{eq:domain-energy} to this resolvent:
\[
 \norm{(R_{i\theta}h)_F}_{\OO_\sigma}^2
 \le C_a\norm{R_{i\theta}h}_{\HH_\sigma}^2+
              C_a\norm h_{\HH_\sigma}\norm{R_{i\theta}h}_{\HH_\sigma}.
\]
Its outgoing trace tends to zero as well. Approximate $g_c$ in
$\HH_\sigma$ by compact smooth $g_j$ supported away from the floor.
For fixed $j$, $Bg_j$ is an ordinary zero-mass datum, so the preceding
argument applies. The uniform completed trace estimate controls
the approximation error independently of $\theta$, proving
\eqref{eq:response-tail}.

Parseval's identity and \eqref{eq:triple-bound} give the stated
periodic bound first on finite Fourier sums, then by completion.
The completed weak equation identifies the limiting derivative
and traces, preserves mass zero, and retains the natural floor input.
The density belongs to $C_{\per}(0,T;\HH_\sigma)$.
Finally the scalar multipliers tend to zero by
\eqref{eq:response-tail}, so finite Fourier truncations approximate
the linear response in operator norm on the same $H^1$ space.
\end{proof}

Keep $a$ fixed and allow $c,m$ near $c_*,m_*$. For a small real
$n\in H^1_{\per}(0,T)$ solve the completed equation
\begin{equation}\label{eq:nonlinear-response}
 h_n=R_TB\{n g_c+n[-c h_n+c^2(h_n)_y]\},\qquad \int_Dh_n(t)=0.
\end{equation}
Here and below $R_TB$ denotes the density component when the entire
triple is understood. The one-dimensional Sobolev product inequality
and Lemma~\ref{lem:periodic-inverse} bound the last term on
$H^1_t\EE_\sigma$ by $C\norm n_{H^1}$ times the triple norm.
A Neumann series therefore gives a unique small analytic solution.
Parameter analyticity follows uniformly in Fourier frequency from
\eqref{eq:parameter-inverse}. Put $F_n=M+h_n$. Its weak equation
has precisely the prescribed mean $m+cn$ and diffusion $a+c^2n$,
with the voltage traces and the natural floor flux. At this stage
$F_n$ is allowed to be signed.

The actual self-consistency equation is the scalar equation
\begin{equation}\label{eq:rate-equation}
 \mathcal F(n,c,m):=n-\ell h_n=0 .
\end{equation}
Its linearization at $n=0$ is the Fourier multiplier
\[
 d_k=1-\resp(ik\Omega_*;m,a,c).
\]
At the base point only $k=\pm1$ are singular. Its restriction to
the complement is boundedly invertible: finitely many nonzero
$d_k$ are bounded away from zero, and
\eqref{eq:response-tail} gives $|d_k|\ge1/2$ in the remaining tail.

\begin{proposition}[A nonconstant self-consistent rate]\label{prop:rate-branch}
For all sufficiently small $r>0$ there are real analytic parameter
values $c(r^2),m(r^2)$ near $c_*,m_*$ and real
$n_r\in H^1_{\per}(0,T)$ solving \eqref{eq:rate-equation}, with exact
first Fourier coefficient $(n_r)_1=r$ and
$\norm{n_r}_{H^1}\to0$.
\end{proposition}

\begin{proof}
Write
$n=ze^{i\Omega_*t}+\overline z e^{-i\Omega_*t}+w$, with
$w_{\pm1}=0$. The complementary equation has an analytic solution
$w=w(z,\overline z,c,m)$ by its bounded inverse and the local
implicit function argument. Equivalently, subtract its linear part
and use the resulting contraction. Time translations commute with
the equation and its projections, so uniqueness makes $w$
translation covariant. In the remaining first coefficient,
a real analytic monomial $z^p\overline z^{\,q}$ must have $p-q=1$.
Consequently
\begin{equation}\label{eq:reduced-equation}
 \mathcal F_1=z\,g(|z|^2,c,m),\qquad
 g(0,c,m)=1-\resp(i\Omega_*;m,a,c).
\end{equation}
The latter equality holds also for nearby parameters: the linear
rate operator remains diagonal in Fourier modes, so its first-mode
input has zero complementary projection. The persistent
complementary inverse thus gives $w_z(0)=0$.

Let $\lambda(c,m)$ be the simple characteristic root near
$i\Omega_*$. Along its crossing curve,
\[
 \lambda_m+\lambda_c\,\partial_m\cH=i\,\partial_m\OmH,\qquad
 \det_{\R}(\lambda_c,\lambda_m)
       =\operatorname{Re}\lambda_c\,\partial_m\OmH\ne0.
\]
The determinant uses the real and imaginary parts as its two rows.
If $d(\lambda,c,m)=1-\resp(\lambda;m,a,c)$, differentiation of
$d(\lambda(c,m),c,m)=0$ multiplies these two parameter columns
by the same nonzero complex number $-d_\lambda$. It therefore
preserves nonsingularity. The real Jacobian of $g(0,c,m)$ in
$(c,m)$ is invertible. The finite-dimensional analytic implicit
function argument now gives $c(\rho),m(\rho)$ with
$g(\rho,c(\rho),m(\rho))=0$ for small real $\rho$.
Take $\rho=r^2$, $z=r>0$. This proves the proposition.
The use of two free physical parameters avoids any assumption
about a cubic coefficient in a one-parameter bifurcation.
\end{proof}

\subsection{Positivity and the physical solution}

An $H^1_t\HH_\sigma$ perturbation of $M$ need not be pointwise
nonnegative. The next argument uses a positive period map instead.

\begin{lemma}[Comparison of prescribed evolutions]\label{lem:flow-comparison}
On each fixed finite interval $[0,t_*]$, let $S_n(t,0)$ be the
prescribed positive conservative evolution with mean $m+cn(t)$
and diffusion $a+c^2n(t)$. For small real $n$ and nearby $m,c$,
\begin{equation}\label{eq:flow-comparison}
 \norm{S_n(t_*,0)-S_*(t_*)}_{\HH_\sigma\to\HH_\sigma}
 \le C_{a,t_*}\bigl(|m-m_*|+\norm n_\infty\bigr),
\end{equation}
where $S_*$ is the frozen evolution at $(m_*,a)$.
\end{lemma}

\begin{proof}
For the same initial datum $f$ put $u_n=S_n(\cdot,0)f$.
The homogeneous forward energy is uniform on a small positive
coefficient neighborhood, so
$\norm{u_n}_{L^\infty_t\HH_\sigma}
 +\norm{(u_n)_y}_{L^2_t\HH_\sigma}\le C_{a,t_*}\norm f_{\HH_\sigma}$.
The difference $u_n-S_*f$ solves the base forced equation with
zero initial datum and input
\[
 g=-[(m-m_*)+cn]u_n+c^2n(u_n)_y .
\]
Its $L^2_t\HH_\sigma$ norm is at most the right side of
\eqref{eq:flow-comparison} times $\norm f$. The completed weak
flux identity and prescribed uniqueness identify this difference
with the solution in Lemma~\ref{lem:causal}. Estimate
\eqref{eq:causal-bound} proves the assertion. This uses a causal
forcing estimate, not an assumed forced renormalization formula.
\end{proof}

\begin{proof}[Proof of Theorem~\ref{thm:main}]
Choose the base point in Lemma~\ref{lem:selection}. By the frozen gap,
there is an integer $p$ with
$\norm{S_*(pT)}_{\HH_{\sigma,0}\to\HH_{\sigma,0}}<1/4$.
Keep this integer fixed. For the solution in
Proposition~\ref{prop:rate-branch}, choose $r>0$ so small that
\eqref{eq:flow-comparison} at $pT$ is less than $1/4$.
The prescribed $pT$ map then contracts every mass-zero difference
by a factor less than $1/2$.

Iterate that map from the normalized nonnegative density $M$.
Every iterate is nonnegative and has mass one. The affine
mass-one space in $\HH_\sigma$ is complete, so the iterates
converge to its unique fixed point $F_0^*$. Positivity and mass
one pass to the limit, since their defining sets are closed.
If $P_n=S_n(T,0)$, periodic coefficients and uniqueness give
$S_n(pT,0)=P_n^p$. The operator $P_n$ commutes with this power,
and $P_nF_0^*$ is another mass-one fixed point of it.
Uniqueness gives $P_nF_0^*=F_0^*$.
Thus $F^*(t)=S_n(t,0)F_0^*$ is nonnegative and $T$-periodic.

The signed response $F_n=M+h_n$ in \eqref{eq:nonlinear-response}
lies in $C_t\HH_\sigma$ with its derivative and traces in the
prescribed energy class. Its completed weak equation is the
same prescribed homogeneous equation. Prescribed uniqueness
identifies it with its evolution from $F_n(0)$. Since it is
$T$-periodic and of mass one, $F_n(0)$ is a fixed point of the
$pT$ map on the full affine space, and hence equals $F_0^*$.
This proves $F_n=F^*\ge0$ without a lower bound for $M$.

Set, with $m=m(r^2)$ and $c=c(r^2)$,
\begin{equation}\label{eq:physical-reconstruction}
 y_*=m-cN_0(m,a),\qquad a_*=a-c^2N_0(m,a).
\end{equation}
At the base point $c=O(a)$ and $N_0$ is uniformly bounded, while
$m>1$. First choose $a$ small and then $r$ small. This makes
$y_*,a_*,c$ positive. The exact rate equation gives
$N=N_0+n_r$, so its first Fourier coefficient is $r>0$.
Moreover
\[
 y_*+cN=m+cn_r,\qquad a_*+c^2N=a+c^2n_r .
\]
The prescribed equation for $F_n$ is therefore the original
nonlinear equation with both feedback terms. The choices
$v_E=2,v_F=1,y_L=1$ satisfy all remaining physical parameter
inequalities.

We verify the requested solution properties. The density belongs
to $H^1_{\per}(0,T;\HH_\sigma)$, hence to
$C_{\per}(0,T;\HH_\sigma)\subset C_{\per}(0,T;L^1(D))$.
The continuous representative has mass one at every time.
Exponential-weight Cauchy--Schwarz bounds all finite conductance
moments uniformly in time, in particular $\int y^2F$.
The positive entropy part is bounded by $\int(F^2+F)$.
For the negative part choose the normalized density
$q_0(v,y)=C e^{-y^2}$. The pointwise inequality
$s\log(s/q_0)\ge s-q_0$, together with the second moment and the
finite positive part, bounds $\int|F\log F|$ uniformly over
the period, with $0\log0=0$.

The low-conductance zero traces and high-conductance flux reset
are the strong completed voltage traces in
\eqref{eq:nonlinear-response}, equivalently those of the
prescribed evolution. Its floor trace is zero for
\[
 Q=(m+cn_r-y)F-(a+c^2n_r)F_y
      =(y_*+cN-y)F-(a_*+c^2N)F_y.
\]
Finally pair $Q$ with any smooth periodic time test and any
smooth test on the closed voltage interval. This pairing lies
in $L^1(\dd y)$ by the weighted energy bounds. The PDE,
integration by parts in $t,v$, and the integrable outgoing and
reset traces put its distributional $y$ derivative in $L^1$
on every tail. It has an absolutely continuous tail
representative with a finite limit at infinity. Its own
$L^1$ integrability forces that limit to be zero. This is the
vanishing distributional flux trace in the theorem.
All required boundary, entropy, continuity, and moment
properties follow for the nonnegative full density $F=F_n$.
\end{proof}
